\section{Hill Cipher}
\label{sec:hill}

\subsection{Konsep}

Hill cipher dikembangkan oleh matematikawan Amerika Lester S. Hill pada tahun 1929. Cipher ini termasuk \textit{cipher} poligram yang berbasis aljabar linear: plainteks dibagi menjadi blok berukuran $m$ huruf, lalu setiap blok dienkripsi dengan $m$ buah persamaan linear modulo 26 \cite{munir_slides_itb,stallings2017}. Untuk $m = 3$ (enkripsi setiap 3 huruf):
\begin{align*}
C_1 &= (k_{11} p_1 + k_{12} p_2 + k_{13} p_3) \bmod 26 \\
C_2 &= (k_{21} p_1 + k_{22} p_2 + k_{23} p_3) \bmod 26 \\
C_3 &= (k_{31} p_1 + k_{32} p_2 + k_{33} p_3) \bmod 26
\end{align*}
atau dalam bentuk matriks
\[
\begin{pmatrix} C_1 \\ C_2 \\ C_3 \end{pmatrix} =
\begin{pmatrix} k_{11} & k_{12} & k_{13} \\ k_{21} & k_{22} & k_{23} \\ k_{31} & k_{32} & k_{33} \end{pmatrix}
\begin{pmatrix} p_1 \\ p_2 \\ p_3 \end{pmatrix} \bmod 26,
\qquad \text{ringkasnya} \quad \mathbf{C} = \mathbf{K}\mathbf{P} \bmod 26.
\]
Matriks $\mathbf{K}$ berukuran $m \times m$ adalah kuncinya.

\begin{contoh}
Kunci
\[
\mathbf{K} = \begin{pmatrix} 17 & 17 & 5 \\ 21 & 18 & 21 \\ 2 & 2 & 19 \end{pmatrix},
\]
plainteks \texttt{paymoremoney}. Tiga huruf pertama, \texttt{pay} $= (15, 0, 24)$, dienkripsi menjadi
\[
\mathbf{C} = \begin{pmatrix} 17 & 17 & 5 \\ 21 & 18 & 21 \\ 2 & 2 & 19 \end{pmatrix}
\begin{pmatrix} 15 \\ 0 \\ 24 \end{pmatrix}
= \begin{pmatrix} 375 \\ 819 \\ 486 \end{pmatrix} \bmod 26
= \begin{pmatrix} 11 \\ 13 \\ 18 \end{pmatrix} = \texttt{LNS}.
\]
Dengan cara yang sama untuk blok \texttt{mor}, \texttt{emo}, dan \texttt{ney}, cipherteks selengkapnya adalah \texttt{LNSHDLEWMTRW}.
\end{contoh}

\subsection{Dekripsi dan Matriks Balikan}

Dekripsi memerlukan matriks balikan $\mathbf{K}^{-1}$ modulo 26 sedemikian sehingga $\mathbf{K}\mathbf{K}^{-1} \equiv \mathbf{I} \pmod{26}$ ($\mathbf{I}$ matriks identitas):
\[
\mathbf{P} = \mathbf{K}^{-1}\mathbf{C} \bmod 26.
\]
$\mathbf{K}^{-1}$ hanya ada jika $\det(\mathbf{K})$ relatif prima dengan 26. Untuk matriks $2 \times 2$:
\[
\mathbf{K} = \begin{pmatrix} a & b \\ c & d \end{pmatrix}
\quad \Rightarrow \quad
\mathbf{K}^{-1} = \det(\mathbf{K})^{-1} \begin{pmatrix} d & -b \\ -c & a \end{pmatrix} \bmod 26,
\qquad \det(\mathbf{K}) = ad - bc,
\]
dengan $\det(\mathbf{K})^{-1}$ balikan modulo 26 dari determinan (Bab~III).

\begin{contoh}
$\mathbf{K} = \begin{pmatrix} 3 & 10 \\ 15 & 9 \end{pmatrix}$, sehingga $\det(\mathbf{K}) = (3)(9) - (10)(15) = -123 \equiv 7 \pmod{26}$. Karena $7^{-1} \equiv 15 \pmod{26}$ (sebab $7 \cdot 15 = 105 \equiv 1$), $-10 \equiv 16$, dan $-15 \equiv 11$:
\[
\mathbf{K}^{-1} = 15 \begin{pmatrix} 9 & 16 \\ 11 & 3 \end{pmatrix}
= \begin{pmatrix} 135 & 240 \\ 165 & 45 \end{pmatrix} \bmod 26
= \begin{pmatrix} 5 & 6 \\ 9 & 19 \end{pmatrix}.
\]
Periksa: $\mathbf{K}\mathbf{K}^{-1} = \begin{pmatrix} 105 & 208 \\ 156 & 261 \end{pmatrix} \bmod 26 = \begin{pmatrix} 1 & 0 \\ 0 & 1 \end{pmatrix}$.

Dengan kunci ini, \texttt{kripto} $= (10, 17), (8, 15), (19, 14)$ dienkripsi blok demi blok. Untuk blok pertama:
\[
\begin{pmatrix} 3 & 10 \\ 15 & 9 \end{pmatrix}\begin{pmatrix} 10 \\ 17 \end{pmatrix}
= \begin{pmatrix} 200 \\ 303 \end{pmatrix} \bmod 26 = \begin{pmatrix} 18 \\ 17 \end{pmatrix} = \texttt{SR},
\]
dan cipherteks lengkapnya \texttt{SRSVPV}. Mengalikan setiap blok cipherteks dengan $\mathbf{K}^{-1}$ mengembalikan \texttt{kripto}.
\end{contoh}

Sebaliknya, matriks $\begin{pmatrix} 6 & 24 \\ 13 & 16 \end{pmatrix}$ tidak dapat dipakai sebagai kunci: determinannya $96 - 312 = -216 \equiv 18 \pmod{26}$, dan $\gcd(18, 26) = 2 \neq 1$, sehingga balikannya tidak ada.

Untuk matriks $3 \times 3$, balikan dapat dicari dengan eliminasi Gauss--Jordan atau dengan rumus adjoin:
\[
\mathbf{K} = \begin{pmatrix} a & b & c \\ d & e & f \\ g & h & i \end{pmatrix}
\quad \Rightarrow \quad
\mathbf{K}^{-1} = \det(\mathbf{K})^{-1} \begin{pmatrix} A & D & G \\ B & E & H \\ C & F & I \end{pmatrix} \bmod 26,
\]
dengan kofaktor
\begin{align*}
A &= ei - fh, & B &= -(di - fg), & C &= dh - eg, \\
D &= -(bi - ch), & E &= ai - cg, & F &= -(ah - bg), \\
G &= bf - ce, & H &= -(af - cd), & I &= ae - bd,
\end{align*}
dan $\det(\mathbf{K}) = aA + bB + cC$.

\begin{contoh}
Untuk kunci $3 \times 3$ pada contoh \texttt{paymoremoney}, $\det(\mathbf{K}) = -939 \equiv 23 \pmod{26}$ dan
\[
\mathbf{K}^{-1} = \begin{pmatrix} 4 & 9 & 15 \\ 15 & 17 & 6 \\ 24 & 0 & 17 \end{pmatrix}.
\]
Dekripsi cipherteks \texttt{LNS} $= (11, 13, 18)$:
\[
\mathbf{P} = \begin{pmatrix} 4 & 9 & 15 \\ 15 & 17 & 6 \\ 24 & 0 & 17 \end{pmatrix}
\begin{pmatrix} 11 \\ 13 \\ 18 \end{pmatrix}
= \begin{pmatrix} 431 \\ 494 \\ 570 \end{pmatrix} \bmod 26
= \begin{pmatrix} 15 \\ 0 \\ 24 \end{pmatrix} = \texttt{pay}.
\]
\end{contoh}

Kekuatan Hill cipher terletak pada penyembunyian frekuensi huruf tunggal: setiap huruf cipherteks bergantung pada semua huruf dalam bloknya, sehingga huruf plainteks yang sama belum tentu dienkripsi menjadi huruf cipherteks yang sama.

\begin{catatan}
\textbf{Kakas daring.} Hill cipher (termasuk perhitungan matriks balikan) dapat dicoba di \url{https://www.dcode.fr/hill-cipher}.
\end{catatan}

\subsection{Implementasi dalam C}

Kode~\ref{lst:hill-c} mengimplementasikan Hill cipher $2 \times 2$. Fungsi \texttt{matrix\_inverse} menolak kunci yang determinannya tidak relatif prima dengan 26.

\begin{ccode}[caption={Hill cipher $2 \times 2$ dalam C}, label={lst:hill-c}]
#include <stdio.h>

int mod26(int x) { return ((x % 26) + 26) % 26; }

int mod_inverse(int a, int m) {       /* algoritma Euclidean diperluas */
    int r0 = m, r1 = ((a % m) + m) % m, t0 = 0, t1 = 1;
    while (r1 != 0) {
        int q = r0 / r1, tmp;
        tmp = r0 - q * r1;  r0 = r1;  r1 = tmp;
        tmp = t0 - q * t1;  t0 = t1;  t1 = tmp;
    }
    if (r0 != 1) return -1;            /* tidak punya invers */
    return ((t0 % m) + m) % m;
}

int matrix_inverse(int K[2][2], int inv[2][2]) {
    int det = mod26(K[0][0] * K[1][1] - K[0][1] * K[1][0]);
    int d = mod_inverse(det, 26);
    if (d < 0) return -1;              /* det tidak relatif prima dengan 26 */
    inv[0][0] = mod26( d * K[1][1]);
    inv[0][1] = mod26(-d * K[0][1]);
    inv[1][0] = mod26(-d * K[1][0]);
    inv[1][1] = mod26( d * K[0][0]);
    return 0;
}

void hill_apply(int M[2][2], const int in[2], int out[2]) {
    out[0] = mod26(M[0][0] * in[0] + M[0][1] * in[1]);
    out[1] = mod26(M[1][0] * in[0] + M[1][1] * in[1]);
}

int main(void) {
    int K[2][2] = {{3, 10}, {15, 9}}, Kinv[2][2];
    int p[2] = {10, 17}, c[2], back[2];   /* "kr" */
    if (matrix_inverse(K, Kinv) != 0) return 1;
    hill_apply(K, p, c);                  /* (18, 17) = "SR" */
    hill_apply(Kinv, c, back);            /* (10, 17) = "kr" */
    printf("%c%c -> %c%c\n", 'A' + back[0], 'A' + back[1], 'A' + c[0], 'A' + c[1]);
    return 0;
}
\end{ccode}

\subsection{Pengayaan: Kriptanalisis Hill Cipher}

Hill cipher mudah dipecahkan dengan \crypt{known-plaintext attack}. Untuk Hill cipher dengan $m = 2$, misalkan kriptanalis mengetahui dua pasangan plainteks--cipherteks:
\[
\mathbf{P} = (19, 7) \rightarrow \mathbf{C} = (0, 23), \qquad \mathbf{P} = (4, 17) \rightarrow \mathbf{C} = (12, 6).
\]
Jadi $\mathbf{K}(19, 7)^T = (0, 23)^T$ dan $\mathbf{K}(4, 17)^T = (12, 6)^T$. Dengan menyusun vektor-vektor tersebut sebagai kolom matriks,
\[
\begin{pmatrix} 0 & 12 \\ 23 & 6 \end{pmatrix} = \mathbf{K} \begin{pmatrix} 19 & 4 \\ 7 & 17 \end{pmatrix}
\quad \Rightarrow \quad \mathbf{C} = \mathbf{K}\mathbf{P} \quad \Rightarrow \quad \mathbf{K} = \mathbf{C}\mathbf{P}^{-1} \bmod 26.
\]
Matriks balikan dari $\mathbf{P}$ adalah
\[
\mathbf{P}^{-1} = \begin{pmatrix} 19 & 4 \\ 7 & 17 \end{pmatrix}^{-1} \bmod 26 = \begin{pmatrix} 25 & 14 \\ 5 & 5 \end{pmatrix},
\]
sehingga
\[
\mathbf{K} = \begin{pmatrix} 0 & 12 \\ 23 & 6 \end{pmatrix} \begin{pmatrix} 25 & 14 \\ 5 & 5 \end{pmatrix} \bmod 26 = \begin{pmatrix} 8 & 8 \\ 7 & 14 \end{pmatrix}.
\]
Secara umum, untuk blok berukuran $m$ cukup $m$ pasangan plainteks--cipherteks yang matriks plainteksnya mempunyai balikan modulo 26.
